\documentclass[11pt]{amsart}

\usepackage[T1]{fontenc}
\usepackage{lmodern}
\usepackage{microtype}
\usepackage{amsmath,amssymb,booktabs}
\usepackage[hidelinks]{hyperref}

\newtheorem{theorem}{Theorem}
\newtheorem{lemma}[theorem]{Lemma}
\newtheorem{proposition}[theorem]{Proposition}
\newtheorem{corollary}[theorem]{Corollary}

\DeclareMathOperator{\dmg}{dmg}
\DeclareMathOperator{\zdmg}{zdmg}

\title[A claw-free zombie-damage counterexample]
      {A counterexample to Davila's claw-free zombie-damage conjecture}
\date{}

\subjclass[2020]{05C57, 05C69, 68R10}
\keywords{damage number, zombies and survivors, claw-free graph,
pursuit-evasion, computer-assisted proof}

\begin{document}

\begin{abstract}
Davila conjectured that every connected claw-free graph $G$ satisfies
$\zdmg(G)\leq 4\dmg(G)$. We construct a connected claw-free graph on nine
vertices with $(\dmg,\zdmg)=(2,9)$, disproving Conjecture~22 of
\cite{Davila}. A ten-vertex extension has values $(2,10)$ and yields
$c_3\geq 5$ in Davila's notation. An exact exhaustive computation shows
that order nine is minimum and that there are exactly $151$ isomorphism
classes of minimum-order counterexamples.
\end{abstract}

\maketitle

\section{The counterexample}

All graphs are finite, simple, and connected. A graph is claw-free if it
has no induced $K_{1,3}$. Write $N(v)$ and $N[v]$ for the open and closed
neighborhoods of a vertex $v$. In the one-cop damage game, the cop and robber
choose initial vertices in that order, the cop moves first, and both players
may move to a neighbor or pass. Capture occurs when the cop moves onto the
robber. A vertex occupied by the robber is \emph{damaged} if the robber
survives the next cop move and then completes a move, possibly a pass, from
that vertex. The damage number $\dmg(G)$ is the number of distinct damaged
vertices under optimal play. In the zombie damage game, the pursuer must move
at each turn to a neighbor one step closer to the survivor, while the survivor
may move or pass; the corresponding value is $\zdmg(G)$. See
\cite{CoxSanaei,Davila}.

Davila's Conjecture~22 states that
\begin{equation}\label{eq:davila}
  \zdmg(G)\leq 4\dmg(G)
\end{equation}
for every connected claw-free graph $G$~\cite[Conjecture~22]{Davila}.

We use the following local form of Davila's trace argument
\cite[Lemma~12]{Davila}.

\begin{lemma}[Local trace lemma]\label{lem:trace}
Let
\[
  a_{-2},a_{-1},a_0,a_1,\ldots,a_m
\]
be a walk such that $a_{i-1}a_i a_{i+1}$ is the unique geodesic between
$a_{i-1}$ and $a_{i+1}$ for every $-1\leq i\leq m-1$. If the zombie starts
at $a_{-2}$ and the survivor starts at $a_0$, then the survivor can damage
every vertex occurring among $a_0,\ldots,a_m$ by following the walk and
passing once at $a_m$.
\end{lemma}

\begin{proof}
Suppose immediately before a zombie move that the zombie is at $a_{i-2}$ and
the survivor is at $a_i$. The unique-geodesic hypothesis forces the zombie
to move to $a_{i-1}$. If $i<m$, the survivor moves to $a_{i+1}$ and damages
$a_i$. Starting with $i=0$ and iterating gives the claim; when $i=m$, the
survivor passes after the zombie moves to $a_{m-1}$, thereby damaging $a_m$.
\end{proof}

For $t\geq1$, define $G_t$ from the cycle
\[
  v_0v_1v_2v_3v_4v_5v_0
\]
by adding a vertex $x$ adjacent to $v_0$ and $v_1$ and a path
\[
  xu_1u_2\cdots u_t.
\]
Thus $|V(G_t)|=7+t$.

\begin{theorem}\label{thm:values}
For $t\in\{2,3\}$, the graph $G_t$ is connected and claw-free, and
\[
  \dmg(G_t)=2,
  \qquad
  \zdmg(G_t)=7+t.
\]
\end{theorem}

\begin{proof}
The only vertices of degree three are $v_0,v_1,x$. Their neighborhoods
contain the edges $v_1x$, $v_0x$, and $v_0v_1$, respectively. Hence $G_t$
is claw-free.

Every vertex is within distance three of $x$, and $d(x,v_3)=3$, so
$\operatorname{ecc}(x)=3$. Also $d(u_t,v_3)=t+3\geq5$. Therefore
\[
  3=\operatorname{ecc}(x)\geq\operatorname{rad}(G_t)
  \geq\left\lceil\frac{\operatorname{diam}(G_t)}2\right\rceil\geq3,
\]
so $\operatorname{rad}(G_t)=3$. The bound
$\dmg(G)\geq\operatorname{rad}(G)-1$ gives $\dmg(G_t)\geq2$
\cite[Theorem~1.3]{HugganMessingerPorter}.

For the reverse inequality, the cop starts at $x$. A robber starting at
$v_0,v_1$, or $u_1$ is captured on the first cop move. If it starts at
$u_j$ with $j\geq2$, the cop moves to $u_1$ and remains there, confining the
robber to $\{u_2,\ldots,u_t\}$, a set of at most two vertices.

For starts on the cycle, use the following two-vertex traps. When the robber
occupies an endpoint of the listed edge, the cop occupies the corresponding
guard vertex.
\[
\begin{array}{c|cc}
\toprule
\text{robber edge} & \text{guard at first endpoint} &
                      \text{guard at second endpoint}\\
\midrule
v_2v_3 & v_0 & v_5\\
v_3v_4 & v_1 & v_0\\
v_4v_5 & v_2 & v_1\\
\bottomrule
\end{array}
\]
In each row the two guards are adjacent, and every move from the listed edge
through an unlisted neighbor lands next to the current guard. Thus the cop
can change guards when the robber crosses the edge and captures on the next
cop move if the robber leaves it. A start at $v_2$ or $v_5$ is placed in the
first or third trap by the move $x\to v_0$ or $x\to v_1$, respectively.
For a start at $v_3$ or $v_4$, the cop passes at $x$ until the robber first
makes a non-pass move and then moves to $v_0$ or $v_1$, respectively. If the
robber never makes such a move, it damages only its initial vertex. Thus at
most two vertices are damaged, proving $\dmg(G_t)\leq2$.

It remains to prove full zombie damage. Write
\[
  U_t=(x,u_1,\ldots,u_t),
\]
and use $U_t$ as a terminal word in the displays below. For the indicated
initial zombie vertices, consider the following walks; in each row the
survivor starts at the third displayed vertex:
\begin{align*}
T_{v_0}&=(v_0,v_5,v_4,v_3,v_2,v_1,
          v_0,v_5,v_4,v_3,v_2,v_1,U_t),\\
T_{v_2}&=(v_2,v_1,v_0,v_5,v_4,v_3,v_2,v_1,U_t),\\
T_{v_3}&=(v_3,v_2,v_1,v_0,v_5,v_4,v_3,v_2,v_1,U_t),\\
T_x&=(x,v_1,v_2,v_3,v_4,v_5,v_0,
      v_1,v_2,v_3,v_4,v_5,v_0,U_t),\\
T_{u_1}&=(u_1,x,v_0,v_5,v_4,v_3,v_2,v_1,U_t).
\end{align*}
The survivor portion of every walk visits all vertices of $G_t$. Every
length-two segment lying on the $6$-cycle or the pendant path is the unique
geodesic between its endpoints. The remaining consecutive triples are, up
to reversal,
\[
  (v_2,v_1,x),\quad (v_5,v_0,x),\quad
  (v_1,x,u_1),\quad (v_0,x,u_1).
\]
For each of these four triples, the endpoints are nonadjacent and the
displayed middle vertex is their unique common neighbor. Hence every
consecutive triple satisfies the hypothesis of Lemma~\ref{lem:trace}.

The automorphism interchanging
\[
  v_0\leftrightarrow v_1,\qquad
  v_2\leftrightarrow v_5,\qquad
  v_3\leftrightarrow v_4
\]
and fixing $x,u_1,\ldots,u_t$ covers initial zombie vertices
$v_1,v_5,v_4$. Finally, if the zombie starts at $u_j$ with $j\geq2$, the
survivor starts at $v_0$ and passes while the zombie is forced toward $u_1$.
After the zombie first reaches $u_1$, the survivor passes once more; the next
zombie turn then begins in the initial position of $T_{u_1}$. Thus the
survivor damages every vertex against every initial zombie position, and
$\zdmg(G_t)=|V(G_t)|=7+t$.
\end{proof}

Davila defines
\[
  c_3=\sup\left\{
  \frac{\zdmg(G)}{\dmg(G)}:
  G\text{ is connected and claw-free},\ \dmg(G)>0
  \right\}.
\]

\begin{corollary}\label{cor:main}
The graph $G_2$ disproves Davila's Conjecture~22, since
\[
  \zdmg(G_2)=9>8=4\dmg(G_2).
\]
Moreover,
\[
  \frac{\zdmg(G_3)}{\dmg(G_3)}=5,
\]
so $c_3\geq5$.
\end{corollary}

\section{Minimum order}

\begin{proposition}\label{prop:minimality}
No connected claw-free graph of order at most eight violates
\eqref{eq:davila}. At order nine, exactly $151$ isomorphism classes violate
\eqref{eq:davila}, all with $(\dmg,\zdmg)=(2,9)$. Hence $G_2$ has minimum
possible order.
\end{proposition}

\begin{proof}[Computational proof]
Let $P\in\{C,Z\}$ denote an ordinary cop or a zombie, write
$\operatorname{val}_C(G)=\dmg(G)$ and
$\operatorname{val}_Z(G)=\zdmg(G)$, and put
\[
  M_C(p,s)=N[p],
  \qquad
  M_Z(p,s)=\{p'\in N(p):d(p',s)=d(p,s)-1\}.
\]
Fix an integer $k$. A state $(D,p,s)$ consists of the damaged set $D$ and
distinct current pursuer and survivor positions $p,s$, immediately before a
pursuer move. Let $W_k^P$ be the least set of states that contains every
state with $|D|\geq k$ and satisfies the following closure rule: add
$(D,p,s)$ if, for every $p'\in M_P(p,s)$, one has $p'\neq s$ and either
$|D\cup\{s\}|\geq k$ or there exists
\[
  s'\in N[s]\setminus\{p'\}
\]
with $(D\cup\{s\},p',s')\in W_k^P$.

A state added after $r$ iterations is one from which the survivor can force
$k$ damaged vertices within at most $r$ further rounds. Conversely, if a
state is outside the stabilized set, failure of the closure rule gives the
pursuer a move that either captures immediately or leaves every legal
survivor reply outside the set while the damaged total remains below $k$.
Repeating such choices prevents the threshold forever. Therefore
\[
  \operatorname{val}_P(G)\geq k
  \quad\Longleftrightarrow\quad
  \forall p\in V(G)\ \exists s\in V(G)\setminus\{p\}:
  (\varnothing,p,s)\in W_k^P.
\]

We evaluated this recurrence exactly on one representative of each connected
isomorphism class through order nine, generated by \texttt{geng -c n} from
\textsc{nauty}~\cite{McKayPiperno}, and retained the claw-free graphs by
checking every triple in each open neighborhood. The results were
\[
\begin{array}{rrrr}
\toprule
n & \text{connected classes} & \text{claw-free classes} &
    \text{counterexample classes}\\
\midrule
1 & 1 & 1 & 0\\
2 & 1 & 1 & 0\\
3 & 2 & 2 & 0\\
4 & 6 & 5 & 0\\
5 & 21 & 14 & 0\\
6 & 112 & 50 & 0\\
7 & 853 & 191 & 0\\
8 & 11117 & 881 & 0\\
9 & 261080 & 4494 & 151\\
\bottomrule
\end{array}
\]
All $151$ counterexamples had values $(2,9)$. Using the vertex order
\[
  v_0,\ldots,v_5,x,u_1,\ldots,u_t,
\]
the graph6 strings of $G_2$ and $G_3$ are, respectively,
\[
  \texttt{HhEM?C@}
  \qquad\text{and}\qquad
  \texttt{IhEM?C@?G}.
\]
\end{proof}

\begin{thebibliography}{9}

\bibitem{CoxSanaei}
D. Cox and A. Sanaei,
\emph{The damage number of a graph},
Australas. J. Combin. \textbf{75} (2019), no.~1, 1--16.

\bibitem{Davila}
R. Davila,
\emph{The Zombie Damage Number of a Graph},
arXiv:2607.16382v1 [math.CO], 2026.

\bibitem{HugganMessingerPorter}
M. A. Huggan, M. E. Messinger, and A. Porter,
\emph{The damage number of the Cartesian product of graphs},
Australas. J. Combin. \textbf{88} (2024), no.~3, 362--384.

\bibitem{McKayPiperno}
B. D. McKay and A. Piperno,
\emph{Practical graph isomorphism, II},
J. Symbolic Comput. \textbf{60} (2014), 94--112.

\end{thebibliography}

\end{document}